%----------------------------------------------------------------------------------------
% Tailored CV — Wissenschaftliche*r Mitarbeiter*in, KI- und Datenethik
% Institut für Medizin- und Datenethik (IMDE), Universität Heidelberg
% ONCOverse Project — Job-ID V000015525
%----------------------------------------------------------------------------------------

\documentclass[10pt,a4paper,sans]{moderncv}

\usepackage[utf8]{inputenc}
\usepackage[T1]{fontenc}

\moderncvstyle{classic}
\moderncvcolor{blue}

\usepackage[scale=0.78]{geometry}
\usepackage{ragged2e}

%----------------------------------------------------------------------------------------

\firstname{Kundai Farai}
\familyname{Sachikonye}

\address{Auf der Lichtung 19}{82178 Puchheim, Germany}
\mobile{(+49) 1626386344}
\email{kundai.sachikonye@bitspark.com}
\homepage{github.com/fullscreen-triangle}{github.com/fullscreen-triangle}

%----------------------------------------------------------------------------------------

\begin{document}

\makecvtitle

%----------------------------------------------------------------------------------------
%	EDUCATION
%----------------------------------------------------------------------------------------

\section{Education}

\cventry{2019--2021}{PhD candidate, Computational Lipidomics}{Technische Universität München / Bayerisches Forschungsinstitut}{}{}{
  ML for large-scale clinical omics cohorts; multi-omics statistical modelling;
  distributed inference across federated clinical sites.
  Departed to pursue independent research.
}
\cventry{2018--2019}{MSc Computational Biology}{Universität Tübingen}{}{}{
  Probabilistic graphical models, statistical bioinformatics, Bayesian methods,
  computational systems biology, ML for biological sequences.
}
\cventry{2014--2017}{MSc Pharmaceutical Biotechnology}{HAW Hamburg}{}{}{}
\cventry{2011--2014}{BSc Biochemistry and Cell Biology}{Jacobs University Bremen}{}{}{}

%----------------------------------------------------------------------------------------
%	RELEVANT MANUSCRIPTS
%----------------------------------------------------------------------------------------

\section{Relevant Manuscripts}

\cvitem{2025}{
  \textbf{On the Thermodynamic Consequences of Categorical Completion on Biological Systems.}
  First-principles AI diagnostic framework with zero trained parameters; MHC-I/II
  binding AUC 0.81 vs.\ NetMHCpan 0.68 ($n=2{,}847$ IEDB peptides);
  neoantigen identification; cancer immune escape predictor.
  \textit{Manuscript. AIMe Registry for Artificial Intelligence.}
}
\cvitem{2025}{
  \textbf{On Disease Epistemology in Blind Receiver.}
  Normative disease state representation $\mathbf{D}\in[0,1]^8$;
  AUC$=1.00$ across 25/25 validation experiments;
  sparse $\ell_1$ LP for minimum-intervention therapeutic target selection;
  first-principles derivation of disease axes (immune, barrier, metabolic,
  circadian). \textit{Manuscript. AIMe Registry / TUM Weihenstephan.}
}
\cvitem{2025}{
  \textbf{On the Epistemology of Purpose in Artificial Intelligence.}
  Formal normative framework for AI system design: VC dimension and
  PAC-learnability bounds; twelve falsifiable predictions; first-principles
  derivation of AI transparency requirements.
  \textit{Manuscript. AIMe Registry for Artificial Intelligence.}
}

%----------------------------------------------------------------------------------------
%	AI ETHICS AND ONCOLOGY TOOLS
%----------------------------------------------------------------------------------------

\section{AI Ethics Foundations: Oncology Systems Built}

\cvitem{syndrome}{
  \textbf{Auditable AI for Cancer Immunopeptidology.}
  3-channel DFT (hydropathy, volume, charge) $\to$ cosine similarity $\rho$;
  zero free parameters trained on patient data; all predictions reproducible
  from algorithm alone, independent of training cohort.
  \textbf{Oncology applications:} neoantigen identification for pancreatic
  cancer (KRAS-mutant PDAC immune evasion); MHC binding prediction;
  escape predictor for therapeutic vaccine design;
  disease state $\mathbf{D}\in[0,1]^8$ with $D_\mathrm{Imm}$ (immune activation)
  and $D_\mathrm{Bar}$ (barrier integrity).
  \textbf{Access design:} browser-deployable (WebGPU), 50\,MB footprint;
  no GPU cluster required.
  Live: \href{https://syndrome-seven.vercel.app/tools}{syndrome-seven.vercel.app/tools}.
}

\cvitem{crown-prince}{
  \textbf{Pharmacological AI with Interpretable Geometry.}
  PD/PK/ADME trajectories derived from molecular geometry; $K_d$ from
  bounded phase-space geometry; Hill equation as limiting case;
  validated 39/39 against NIST with zero empirical parameters.
  Ethical property: therapeutic target predictions are geometrically
  grounded, not statistical associations from potentially biased training data.
  \href{https://github.com/fullscreen-triangle/crown-prince}{github.com/fullscreen-triangle/crown-prince}
}

\cvitem{gospel}{
  \textbf{Federated Multi-Omics Infrastructure.}
  $10^6$ variants/second; Bayesian network multi-omics integration
  (host genotype $\to$ metagenomics $\to$ transcriptomics conditional graph);
  validated: 1000~Genomes, gnomAD~v3.1.2, UK~Biobank.
  Relevant to ONCOverse data governance: practical experience of what data flows
  across multi-site clinical studies, where pseudonymisation fails,
  and what consent architecture the analysis methods require.
  \href{https://github.com/fullscreen-triangle/gospel}{github.com/fullscreen-triangle/gospel}
}

%----------------------------------------------------------------------------------------
%	AI ETHICS EXPERTISE
%----------------------------------------------------------------------------------------

\section{AI Ethics Competencies}

\cvitem{Explainability}{
  Technical specification of algorithmic transparency beyond procedural compliance:
  distinction between post-hoc explainability (SHAP, LIME) and structural
  auditability (first-principles, zero-parameter models); analysis of when each
  is ethically sufficient for clinical decision support
}

\cvitem{Algorithmic Fairness}{
  Identification of demographic bias entry points in clinical AI pipelines:
  training data composition, threshold selection, population generalisability
  gaps, validation cohort vs.\ deployment context mismatch; quantitative
  fairness metrics (equalized odds, calibration across subgroups)
}

\cvitem{Access Equity}{
  Infrastructure-level access analysis: delivery architecture as determinant
  of equitable access (browser vs.\ GPU cluster); computational cost as
  barrier to adoption in resource-limited clinical settings;
  design-time vs.\ post-deployment access intervention
}

\cvitem{Data Governance}{
  Multi-site clinical data ethics: re-identification risk in pseudonymised
  multi-omics data; consent architecture constraints imposed by analysis methods;
  federated vs.\ centralised data pooling tradeoffs; GDPR and clinical trial
  data protection in AI research contexts
}

\cvitem{Normative Frameworks}{
  Formal treatment of normative commitments in AI system design:
  disease ontology as ethical choice; threshold selection as value judgement;
  minimum-intervention principles in therapeutic AI; patient autonomy
  in algorithmic clinical decision support
}

%----------------------------------------------------------------------------------------
%	EXPERIENCE
%----------------------------------------------------------------------------------------

\section{Experience}

\cventry{2021--present}{Independent AI \& Computational Biology Researcher}{Bitspark GmbH}{}{}{
  Design, implementation, validation, and deployment of AI diagnostic systems
  in oncology, pharmacology, immunopeptidology, and multi-omics.
  All systems publicly deployed with independently verifiable validation numbers.
}

\cventry{2019--2021}{Computational Lipidomics}{TUM / Bayerisches Forschungsinstitut}{Munich}{}{
  ML for clinical omics cohorts; distributed pipeline maintenance; multi-omics
  integration across federated clinical sites. Python, Java.
}

\cventry{2017--2018}{Glycomics and Proteomics}{Universitätsklinikum Hamburg-Eppendorf}{}{}{
  Automated LC-MS/MS and MALDI-TOF pipeline; clinical sample processing;
  data annotation and quality control for clinical cohort studies.
}

%----------------------------------------------------------------------------------------
%	LANGUAGES
%----------------------------------------------------------------------------------------

\section{Languages}

\cvitemwithcomment{English}{Native / Fluent}{}
\cvitemwithcomment{German}{Conversational}{Living in Germany continuously since 2011}

%----------------------------------------------------------------------------------------

\renewcommand{\listitemsymbol}{-~}

\end{document}
